\documentclass[10pt]{exam}
\printanswers
\usepackage{epsfig,amssymb}
\usepackage{xcolor}
\definecolor{darkred}{rgb}{0.5,0,0}
\definecolor{darkgreen}{rgb}{0,0.5,0}
\usepackage{hyperref}
\usepackage{fullpage}
\usepackage{tikz}
\pagestyle{empty}
\usepackage{listings}
\setlength{\parindent}{0pt}
\setlength{\parskip}{.25cm}
\newcommand{\comment}[1]{}
\boxedpoints
\addpoints
\usepackage{amsmath}
\usepackage{algorithm2e}
\usepackage{url}
\usepackage{enumitem}
\usepackage{graphics}

\pagestyle{headandfoot}
\firstpageheader{\Large Math 309}{{\Large Homework 3}}{\Large Fall 2026}

\begin{document}

\hrule

\vspace{.50cm} Name:
\text{Charlie Pyle}
\vspace{.25cm}
\begin{questions}

\question Let $A = \{1, 2, \{3,4\}, \{5\}\}$. Decide which of each of the following statements are true or false.

\begin{parts}
\part $1 \in A$
\part $\{1\} \in A$
\part $\{5\} \in A$
\part $\{1,2\} \in A$
\part $\{1,2\} \subseteq A$
\part $\emptyset \in A$
\end{parts}

\begin{solution}
Enter your solution here:\\
\textbf{(a)} $1 \in A$ \quad \textbf{True}. The element $1$ is listed explicitly as a member of $A$.

\textbf{(b)} $\{1\} \in A$ \quad \textbf{False}. The set $\{1\}$ is not one of the four elements listed in $A$; only the element $1$ itself belongs to $A$.

\textbf{(c)} $\{5\} \in A$ \quad \textbf{True}. The set $\{5\}$ is listed explicitly as a member of $A$.

\textbf{(d)} $\{1,2\} \in A$ \quad \textbf{False}. The set $\{1,2\}$, taken as a single object, is not one of the four elements listed in $A$.

\textbf{(e)} $\{1,2\} \subseteq A$ \quad \textbf{True}. Both $1 \in A$ and $2 \in A$, so every element of $\{1,2\}$ is an element of $A$.

\textbf{(f)} $\emptyset \in A$ \quad \textbf{False}. The empty set is not one of the four elements listed in $A$.
\end{solution}

\question Consider the following logical equivalence:

\begin{center}
\textbf{Theorem.} For all logical statements $P$, $Q$, and $R$, we have
\[(P \Rightarrow Q) \Rightarrow R \equiv \bigl((\lnot P) \Rightarrow R\bigr) \land (Q \Rightarrow R).\]
\end{center}

Prove Theorem 1 using previously established logical equivalences.

\begin{solution}
Enter your solution here:\\
\textit{Proof.} We assume that $P$, $Q$, and $R$ are logical statements, and we will show that $(P \Rightarrow Q) \Rightarrow R$ is logically equivalent to $\bigl((\lnot P) \Rightarrow R\bigr) \land (Q \Rightarrow R)$ by simplifying the right-hand side with a chain of logical equivalences until it matches the left-hand side.
\begin{align*}
\bigl((\lnot P) \Rightarrow R\bigr) \land (Q \Rightarrow R)
&\equiv (\lnot \lnot P \lor R) \land (\lnot Q \lor R)
&& \text{Implication (applied twice)} \\
&\equiv (P \lor R) \land (\lnot Q \lor R)
&& \text{Double Negation} \\[2pt]
&\equiv (P \land \lnot Q) \lor R
&& \text{Distributive} \\[2pt]
&\equiv (\lnot \lnot P \land \lnot Q) \lor R
&& \text{Double Negation} \\[2pt]
&\equiv \lnot (\lnot P \lor Q) \lor R
&& \text{De Morgan} \\[2pt]
&\equiv \lnot (P \Rightarrow Q) \lor R
&& \text{Implication} \\[2pt]
&\equiv (P \Rightarrow Q) \Rightarrow R
&& \text{Implication}
\end{align*}

Since each step above is a logical equivalence, we conclude that
\[(P \Rightarrow Q) \Rightarrow R \equiv \bigl((\lnot P) \Rightarrow R\bigr) \land (Q \Rightarrow R).\]
$\blacksquare$
\end{solution}

\question Use set builder notation to write down the following sets:

\begin{parts}
\part The set $A$ of all odd positive integers.
\part The set $B$ of all real numbers strictly larger than $\pi$.
\part The set $C$ given by
\[C = \left\{ \ldots, 2^{-4}, 2^{-3}, 2^{-2}, \frac{1}{2}, 1 \right\}.\]
\end{parts}

\begin{solution}
Enter your solution here:\\
\textbf{(a)}
\begin{align*}
A = \{x \mid x = 2k+1, \ k \in \mathbb{Z}, \ k \geq 0\}
\end{align*}

\textbf{(b)}
\begin{align*}
B = \{x \in \mathbb{R} \mid x > \pi\}
\end{align*}

\textbf{(c)} 
\begin{align*}
C = \{x \mid x = 2^{n}, \ n \in \mathbb{Z}, \ n \leq 0\}
\end{align*}
\end{solution}

\question Consider the following statements:
\begin{parts}
\part There exists a rational number that is less than or equal to $-1010$.
\part For all integers $x$ and $y$, we have $xy > x$.
\part For all real numbers $\epsilon$, there exists a real number $\delta$ such that if $|x - \pi| < \delta$, then $|\sin(x)| < \epsilon$.
\end{parts}

For each statement, complete the following three tasks:
\begin{itemize}
\item Rewrite the statement in symbolic form.
\item Write its negation in symbolic form.
\item Write a useful negation of each of the following statements as English sentences, without symbols for quantifiers, but where you can use symbols as they were used in the original statements.
\end{itemize}

\begin{solution}
Enter your solution here:\\

\textbf{(a)} There exists a rational number that is less than or equal to $-1010$.

Symbolic form:
\begin{align*}
\exists x \in \mathbb{Q} \ (x \leq -1010)
\end{align*}

Negation (symbolic form):
\begin{align*}
\lnot \bigl[\exists x \in \mathbb{Q} \ (x \leq -1010)\bigr] \equiv \forall x \in \mathbb{Q} \ (x > -1010)
\end{align*}

Useful negation (English sentence): Every rational number $x$ satisfies $x > -1010$.

\vspace{.2cm}
\textbf{(b)} For all integers $x$ and $y$, we have $xy > x$.

Symbolic form:
\begin{align*}
\forall x \in \mathbb{Z} \ \forall y \in \mathbb{Z} \ (xy > x)
\end{align*}

Negation (symbolic form):
\begin{align*}
\lnot \bigl[\forall x \in \mathbb{Z} \ \forall y \in \mathbb{Z} \ (xy > x)\bigr] \equiv \exists x \in \mathbb{Z} \ \exists y \in \mathbb{Z} \ (xy \leq x)
\end{align*}

Useful negation (English sentence): There exist integers $x$ and $y$ such that $xy \leq x$.

\vspace{.2cm}
\textbf{(c)} For all real numbers $\epsilon$, there exists a real number $\delta$ such that if $|x - \pi| < \delta$, then $|\sin(x)| < \epsilon$.

Symbolic form:
\begin{align*}
\forall \epsilon \in \mathbb{R} \ \exists \delta \in \mathbb{R} \ \bigl(|x - \pi| < \delta \rightarrow |\sin(x)| < \epsilon\bigr)
\end{align*}

Negation (symbolic form):
\begin{align*}
\lnot \bigl[\forall \epsilon \in \mathbb{R} \ \exists \delta \in \mathbb{R} \ \bigl(|x - \pi| < \delta \rightarrow |\sin(x)| < \epsilon\bigr)\bigr]
\equiv \exists \epsilon \in \mathbb{R} \ \forall \delta \in \mathbb{R} \ \bigl(|x - \pi| < \delta \land |\sin(x)| \geq \epsilon\bigr)
\end{align*}

Useful negation (English sentence): There exists a real number $\epsilon$ such that for every real number $\delta$, we have $|x - \pi| < \delta$ and $|\sin(x)| \geq \epsilon$.

\end{solution}

\end{questions}

\end{document}
